\documentclass{article}
\usepackage[utf8]{inputenc}
\usepackage{geometry}
\geometry{a4paper, margin=1in}
\usepackage{xcolor}
\usepackage{setspace}

\definecolor{poemtitle}{RGB}{236, 0, 140} % Approximate magenta color from the original title

\title{\Huge\color{poemtitle}\fontfamily{pzc}\selectfont Father to Son}
\author{\Large\itshape Elizabeth Jennings}
\date{}

\begin{document}

\maketitle

\begin{center}
\begin{minipage}{0.65\textwidth}
\setstretch{1.3}
\setlength{\parindent}{0pt}
\setlength{\parskip}{1em}

I do not understand this child\\
Though we have lived together now\\
In the same house for years. I know\\
Nothing of him, so try to build\\
Up a relationship from how\\
He was when small. Yet have I killed\\

The seed I spent or sown it where\\
The land is his and none of mine?\\
We speak like strangers, there's no sign\\
Of understanding in the air.\\
This child is built to my design\\
Yet what he loves I cannot share.\\

Silence surrounds us. I would have\\
Him prodigal, returning to\\
His father's house, the home he knew,\\
Rather than see him make and move\\
His world. I would forgive him too,\\
Shaping from sorrow a new love.\\

Father and son, we both must live\\
On the same globe and the same land,\\
He speaks: I cannot understand\\
Myself, why anger grows from grief.\\
We each put out an empty hand,\\
Longing for something to forgive.

\end{minipage}
\end{center}

\end{document}
\documentclass{article}
\usepackage[utf8]{inputenc}
\usepackage{geometry}
\geometry{a4paper, margin=1in}
\usepackage{enumitem}
\usepackage{xcolor}

% Color definition for headers matching the original style
\definecolor{notespink}{RGB}{236, 0, 140} 

\begin{document}

% Header section
\noindent 72 \hfill \textsc{Hornbill}

\vspace{2em}

% Main title section
\noindent {\Large\bfseries Think it out}

\vspace{1em}

% Questions section
\begin{enumerate}[label=\arabic*., leftmargin=2em, itemsep=0.5em]
    \item Does the poem talk of an exclusively personal experience or is it fairly universal?
    \item How is the father's helplessness brought out in the poem?
    \item Identify the phrases and lines that indicate distance between father and son.
    \item Does the poem have a consistent rhyme scheme?
\end{enumerate}

\vspace{3em}

% Notes title line decoration
\begin{center}
    {\color{notespink}\vrule height 0.5ex width 3em}\ \ {\Large\bfseries\color{notespink} Notes}\ \ {\color{notespink}\vrule height 0.5ex width 3em}
\end{center}

\vspace{1.5em}

% Notes content
\noindent The poem is autobiographical in nature and describes the relationship between a father and his son.

\vspace{1.5em}

\noindent {\bfseries Understanding the poem}

\vspace{0.8em}

\noindent Questions are based on
\begin{itemize}[label=\textbullet, leftmargin=2em, itemsep=0.3em]
    \item the universality of the experience described
    \item phrases in the poem
    \item rhyme scheme in the poem
\end{itemize}

\end{document}

